\section*{Capítulo \ref{ch:b1:elimination} --- \texorpdfstring{$\beta$}{Beta}-eliminación}

\begin{solution}{exo:b1:elimination:1}
2-Bromobutano: but-1-eno, \cip{E}- y \cip{Z}-but-2-eno; el mayoritario es el
\cip{E}-but-2-eno. 2-Bromo-2-metilbutano: 2-metilbut-2-eno (mayoritario,
trisustituido) y 2-metilbut-1-eno. Bromociclohexano: solo ciclohexeno.
\end{solution}

\begin{solution}{exo:b1:elimination:2}
E1: $v = k[\mathrm{RBr}]$; E2: $v = k[\mathrm{RBr}][\ce{EtO-}]$. Se mide la
\omterm{def:b1:rate-laws:initial-rate}{velocidad inicial} de formación del alqueno a varias concentraciones de
etóxido: crece en proporción en la E2 y no cambia en la E1.
\end{solution}

\begin{solution}{exo:b1:elimination:3}
A lo largo de C2--C3, el Br de C2 puede estar en anti respecto de
cualquiera de los dos hidrógenos de C3 o del grupo metilo C4: las dos
\omterm{def:b1:stereochemistry-in-depth:isomers}{conformaciones} con un H en anti respecto del Br pueden reaccionar por E2
(dando los alquenos \cip{E} y \cip{Z}); la tercera no puede eliminar
hacia C3.
\end{solution}

\begin{solution}{exo:b1:elimination:4}
SN2 (primario, hidróxido, temperatura baja); E2 (terciario, \omterm{def:b1:predominant-reaction:strong-acid}{base fuerte} y
voluminosa); SN2 (yoduro, buen \omterm{def:b1:electronic-effects:nucleophile}{nucleófilo} y \omterm{def:b1:predominant-reaction:strong-acid}{base débil}, \omterm{def:b1:intermolecular-forces-solvents:solvent-classes}{disolvente
aprótico}); SN1 y E1, con más E1 al calentar.
\end{solution}

\begin{solution}{exo:b1:elimination:5}
C3 lleva un solo hidrógeno. En cada diastereoisómero, situar ese H en
anti respecto del Br de C2 fija las posiciones de los demás grupos: el
metilo de C2 y el etilo de C3 están del mismo lado en un diastereoisómero
y en lados opuestos en el otro. Uno da el \cip{E}-, y el otro, el
\cip{Z}-3-metilpent-2-eno: una \omterm{def:b1:nucleophilic-substitution:stereospecific}{reacción estereoespecífica}.
\end{solution}

\begin{solution}{exo:b1:elimination:6}
El OH se protona, \ce{R-OH2+}; sale el agua y da el \omterm{def:b1:electronic-effects:intermediates}{carbocatión} terciario
\ce{(CH3)2C+-CH2CH3}; una base (agua, \ce{HSO4-}) arranca un protón de un
carbono vecino. Producto mayoritario (Zaitsev): 2-metilbut-2-eno.
\end{solution}

\begin{solution}{exo:b1:elimination:7}
El grupo terc-butilo bloquea la \omterm{def:b1:stereochemistry-in-depth:conformations}{silla} en la que él es \omterm{def:b1:stereochemistry-in-depth:conformations}{ecuatorial}. En el
isómero \textit{trans}, el bromo es entonces también \omterm{def:b1:stereochemistry-in-depth:conformations}{ecuatorial}, nunca en
anti respecto de un C--H: la E2 debe esperar a la \omterm{def:b1:stereochemistry-in-depth:conformations}{silla} invertida, muy
desfavorable. En el isómero \textit{cis}, el bromo es \omterm{def:b1:stereochemistry-in-depth:conformations}{axial}, con dos
hidrógenos trans-diaxiales: E2 rápida.
\end{solution}

\begin{solution}{exo:b1:elimination:8}
El etóxido, pequeño, arranca el hidrógeno que da el alqueno más
sustituido (Zaitsev). El terc-butóxido, voluminoso, alcanza con más
facilidad los hidrógenos expuestos del grupo metilo que el hidrógeno del
\ce{CH2} interno: el alqueno menos sustituido.
\end{solution}

\begin{solution}{exo:b1:elimination:9}
El bromometano no tiene ningún carbono $\beta$. En el
1-bromo-2,2-dimetilpropano, el carbono $\beta$ es cuaternario y no lleva
ningún hidrógeno.
\end{solution}

\begin{solution}{exo:b1:elimination:10}
Ambos productos necesitan el \omterm{def:b1:electronic-effects:intermediates}{carbocatión}, que se forma en la etapa lenta
a la velocidad $k_1[\mathrm{RBr}]$. Después se reparte entre la captura
por el etanol (\omterm{def:b1:rate-laws:order}{constante de velocidad} $k_S$) y la pérdida de un protón
(\omterm{def:b1:rate-laws:order}{constante de velocidad} $k_E$), ambas de pseudoprimer orden respecto del
disolvente: la fracción de alqueno es $k_E/(k_E + k_S)$, independiente de
la concentración del \omterm{def:b1:nucleophilic-substitution:substitution}{sustrato}.
\end{solution}

\begin{solution}{exo:b1:elimination:11}
Partiendo de la \omterm{def:b1:stereochemistry-in-depth:isomers}{conformación} del \omterm{def:b1:stereochemistry-in-depth:enantiomers}{compuesto meso} con los dos Br en anti
(los dos fenilos, entonces, en anti, y los dos H en anti), se gira C1
$120^\circ$ para situar su H en anti respecto del Br de C2. Los dos grupos
fenilo quedan entonces del mismo lado del eje H--Br: terminan en
\textit{cis}. En C1, el Br tiene prioridad sobre el Ph; en C2, el Ph sobre
el H; el Br y el fenilo de C2 están en \textit{trans}:
\cip{E}-1-bromo-1,2-difenileteno.
\end{solution}

\begin{solution}{exo:b1:elimination:12}
Solo reacciona la \omterm{def:b1:stereochemistry-in-depth:conformations}{silla} con el \omterm{def:b1:electronic-effects:nucleophile}{grupo saliente} \omterm{def:b1:stereochemistry-in-depth:conformations}{axial}:
$v = k\,x\,[\mathrm{RY}][\mathrm{B}]$,
con $x$ la fracción de esa \omterm{def:b1:stereochemistry-in-depth:conformations}{silla}. Cada grupo alquilo \omterm{def:b1:stereochemistry-in-depth:conformations}{axial} que exige
cuesta varios kJ/mol (\qty{5.7}{kJ/mol} para un metilo), lo que divide
$x$ por unos diez por grupo a temperatura ambiente: con dos o tres grupos
de este tipo, $x$ baja a $10^{-2}$ o $10^{-3}$, y la reacción es otro
tanto más lenta.
\end{solution}

\begin{solution}{pb:b1:elimination:1}
\textbf{1.} Una \omterm{def:b1:stereochemistry-in-depth:conformations}{silla} con el Cl, el isopropilo y el metilo todos
\omterm{def:b1:stereochemistry-in-depth:conformations}{ecuatoriales}.
\textbf{2.} Isopropilo y metilo \omterm{def:b1:stereochemistry-in-depth:conformations}{ecuatoriales}, Cl \omterm{def:b1:stereochemistry-in-depth:conformations}{axial}.
\textbf{3.} \omterm{def:b1:stereochemistry-in-depth:enantiomers}{Diastereoisómeros}: solo difieren en C1.
\textbf{4.} \omterm{def:b1:electronic-effects:nucleophile}{Grupo saliente} e hidrógeno $\beta$ trans-diaxiales.
\textbf{5.} Solo en la \omterm{def:b1:stereochemistry-in-depth:conformations}{silla} invertida, con el Cl, el isopropilo y el
metilo todos \omterm{def:b1:stereochemistry-in-depth:conformations}{axiales}.
\textbf{6.} Los tres grupos \omterm{def:b1:stereochemistry-in-depth:conformations}{axiales}: solo el metilo cuesta unos
\qty{5.7}{kJ/mol}; el isopropilo, mayor, más, y el cloro, un poco: en
total, unos \qty{13}{kJ/mol}, de modo que una fracción de alrededor de
$\exp(-13000/2479) = \num{5e-3}$, el valor que toma el problema.
\textbf{7.} C2 (un H, pues lleva el grupo isopropilo) y C6 (dos H).
\textbf{8.} Con el Cl \omterm{def:b1:stereochemistry-in-depth:conformations}{axial} en C1, el H \omterm{def:b1:stereochemistry-in-depth:conformations}{axial} de C2 (pues el isopropilo es
\omterm{def:b1:stereochemistry-in-depth:conformations}{ecuatorial}) y el H \omterm{def:b1:stereochemistry-in-depth:conformations}{axial} de C6: dos.
\textbf{9.} En la \omterm{def:b1:stereochemistry-in-depth:conformations}{silla} triaxial, el grupo isopropilo ocupa la posición
\omterm{def:b1:stereochemistry-in-depth:conformations}{axial} de C2, de modo que su H es \omterm{def:b1:stereochemistry-in-depth:conformations}{ecuatorial}: no está en anti. C6 tiene un
H \omterm{def:b1:stereochemistry-in-depth:conformations}{axial}: en anti.
\textbf{10.} El cloruro de neomentilo, dos; el cloruro de mentilo, uno.
\textbf{11.} A lo largo de C1--C6: el Cl (\omterm{def:b1:stereochemistry-in-depth:conformations}{axial}, arriba) en C1, opuesto al
H \omterm{def:b1:stereochemistry-in-depth:conformations}{axial} (abajo) de C6; los enlaces del anillo y el H \omterm{def:b1:stereochemistry-in-depth:conformations}{ecuatorial},
alternados entre ellos.
\textbf{12.} En una \omterm{def:b1:stereochemistry-in-depth:conformations}{silla}, ningún enlace C--H es \omterm{def:b1:elimination:periplanar}{sinperiplanar} respecto
del enlace C--Cl; forzarlo supondría un anillo eclipsado y tensionado.
\textbf{13.} Al arrancar el H de C2: 4-metil-1-(propan-2-il)ciclohex-1-eno
(trisustituido); al arrancar un H de C6:
3-metil-6-(propan-2-il)ciclohex-1-eno (disustituido).
\textbf{14.} El alqueno trisustituido, el más sustituido y estable.
\textbf{15.} Los dos alquenos, con un \qty{78}{\percent} del
trisustituido: coherente.
\textbf{16.} Solo el alqueno disustituido: lo contrario de la predicción
de Zaitsev, impuesto por el único hidrógeno anti disponible.
\textbf{17.} Ambas son estereoespecíficas (exigencia anti); el cloruro de
neomentilo reacciona de forma regioselectiva (78/22); el cloruro de
mentilo da un solo isómero constitucional.
\textbf{18.} El carbono que lleva el Cl es secundario y está flanqueado
por un grupo isopropilo: impedido para la SN2, mientras que el etóxido es
una \omterm{def:b1:predominant-reaction:strong-acid}{base fuerte}.
\textbf{19.} $v \propto x\,n_{\mathrm{H}}$: fracción de \omterm{def:b1:stereochemistry-in-depth:conformations}{silla} reactiva por
número de hidrógenos anti.
\textbf{20.} $0.95 \times 2 = 1.90$.
\textbf{21.} $\num{5.0e-3} \times 1 = \num{5.0e-3}$.
\textbf{22.} \textbf{$k(\text{neomentilo})/k(\text{mentilo}) = 1.90/\num{5.0e-3}
= 380$}.
\end{solution}
